% Chapter 7 -- Redirecting Input and Output.
% Sources (Law 1): os/milton/command.c (cmd_redir_parse, run_with_redirect,
% run_output_and_dispatch: '>' create, '>>' append, '<' input, MSG-DOS-0046
% when the target cannot be created, MSG-DOS-0005 when its writes do not fit;
% cmd_pipe_split and run_pipeline: '|' through AH=5Bh temporary files,
% MSG-DOS-0005 when one cannot be created or filled; initech-66oc,
% test-redir-full-emu); devices.c NUL.  The filters MORE, FIND and SORT are
% built and gated (os/milton/{more,find,sort}_program.asm; test-more-filter,
% test-find-filter, test-sort-filter) and ship in DOS (initech-2kgh/zc68,
% Makefile FLAIR_DOS_EXTERNALS, test-dos-folder).
% GATES: test-redir test-redir-parse test-hsct-redir test-bsy7-redir test-bsy8-pipe test-bsy9-redir test-redir-full-emu test-pubs-screens
\Gates{test-redir, test-redir-parse, test-hsct-redir, test-bsy7-redir, test-bsy8-pipe, test-bsy9-redir, test-redir-full-emu, test-pubs-screens}
\chapter{Redirecting Input and Output}

\section{Standard Input and Standard Output}
\index{standard input}\index{standard output}
Most commands read what they need from the keyboard and write their results
to the display.  The keyboard is their \term{standard input} and the
display their \term{standard output}.  InitechDOS lets you change both for
a single command: you can send what a command displays to a file, or have a
command read from a file what it would otherwise read from the keyboard.
This is called \term{redirection}.

\section{Sending Output to a File}
\index{redirection}
Put \type{>} and a file name after a command to send its output to that file
instead of the display.  If the file exists, it is replaced.  Use \type{>>}
to add the output to the end of an existing file instead.  If the file is
read-only, InitechDOS displays \texttt{File creation error} and does not run
the command; if the output does not fit on the disk, it displays
\texttt{Insufficient disk space}.
\begin{Steps}
\item Type \type{CLS} and press \key{Enter}.
\item Type \type{VER > LIST.TXT} and press \key{Enter}.  Nothing is
  displayed: the output went to \file{LIST.TXT}.
\item Type \type{ECHO Second line >> LIST.TXT} and press \key{Enter}.
\item Type \type{TYPE LIST.TXT} and press \key{Enter}.
\end{Steps}

\Screen[rows=8, label=fig:redir]{screens/redir.txt}{Output sent to a file}

\file{LIST.TXT} holds the output of both commands.  You can redirect the
output of any command this way: \type{DIR > FILES.TXT} saves a directory
listing to print or keep.  To throw output away, send it to the device
\file{NUL}: \type{DIR > NUL}.

\section{Taking Input from a File}
Put \type{<} and a file name after a command to make it read from the file
what it would otherwise read from the keyboard.  For example, \cmd{LABEL}
asks for the new label; here the answer comes from a file made with
\cmd{ECHO}:
\begin{Steps}
\item Type \type{CLS} and press \key{Enter}.
\item Type \type{ECHO REPORTS > NAME.TXT} and press \key{Enter}.
\item Type \type{LABEL < NAME.TXT} and press \key{Enter}.
\item Type \type{VOL} and press \key{Enter}.
\end{Steps}

\Screen[rows=7, label=fig:stdin]{screens/stdin.txt}{Input taken from a file}

\cmd{LABEL} read \texttt{REPORTS} from \file{NAME.TXT}, so its question is
not answered on the screen; \cmd{VOL} shows the new label.  Put the original
label back with \type{LABEL INITECH}.

\section{Pipes and Filters}
\index{pipes}\index{filters}
A \term{pipe}, the character \type{|} between two commands, sends the output
of the first command to the second as its input.  InitechDOS keeps the text
in a temporary file on the disk while it passes it on, under a name no other
file is using, and deletes it afterward; if there is no room for the file,
it displays \texttt{Insufficient disk space} and does not run the second
command.

A \term{filter} is a program that reads its standard input, changes it, and
writes the result to its standard output.  Filters are most useful in
pipes.
InitechDOS provides three:

\SideHead{MORE}Displays its input one screen at a time.  \type{TYPE
PHONE.TXT | MORE}

\SideHead{FIND}Displays the lines that contain a string.  \type{TYPE
PHONE.TXT | FIND "Help"}

\SideHead{SORT}Sorts the lines of its input.  \type{SORT < PHONE.TXT >
SORTED.TXT}
